\section{Terminal passages and their independent arguments}
\label{sec:terminal}

The real subdivision supplies identities for scalar fields.  We now turn
these identities into actual transfer functions to which the packet theorems
apply.  In particular, a normalization involving a proposed endpoint must
remain an identity when that endpoint does not match the computed flow.

\begin{theorem}[Terminal reduction]\label{thm:terminal-reduction}
Fix the finite preparation data, exponent lists, and geometric tolerances of
Theorem~\ref{thm:finite-templates}.  After finitely many further real
subanalytic subdivisions, every terminal passage is represented by one scalar
transfer equation and regular, locally unique argument graphs.  The transfer
is an ordinary analytic transfer, an additive transfer of
Theorem~\ref{thm:additive-packets}, a one-rate transfer of
Theorem~\ref{thm:regular-packets}, or a mixed transfer of
Theorem~\ref{thm:mixed-packets}.

More precisely, the following statements hold.
\begin{enumerate}
\item These are actual real ODE transfers on open independent-argument
domains.  The domains have finite lists of strict clock, sign, divisor,
analytic-box, amplitude, and state-room margins.  A bounded limit at which
these margins remain positive is an analytic point of the transfer family.
At bounded-clock corners, adding the indicated bounded ratios gives analytic
normalized transfers also at zero clock limits, with state room.
\item At fixed physical parameters, on the graph of the argument ties near
an actual passage, the equation is the physical mismatch in regular
transverse coordinates.  This identity holds for independently proposed
nearby endpoints, before imposing the transfer equation.
\item For an arbitrary sequence of arguments and any fixed finite collection
of independent-argument derivatives, the transfers and those derivatives
reduce to the preceding packet primitives and bounded analytic operations.
This assertion includes sequences on which any of the imposed strict
margins tend to zero; it does not assume the positive-margin condition in
the first assertion.  When all enlarged normalized arguments stay bounded,
the kernel factors have analytic extensions with state room after the
finite graph additions, and algebraic divisors can be cleared.  If an added
argument diverges, its recoverable coordinate is included in the packet
test.
The reduction may introduce bounded ratios and a fixed finite internal
subdivision chosen for that sequence.  The old derivatives are obtained by
the chain rule for exact transfer identities.  The internal state and finite
jet graphs are locally unique and preserve isolated solutions.
The coefficient determinations are those fixed by the primitive theorems
and are preserved under their allowed retests.
\end{enumerate}
The number of geometric passage types and boundary crossings is finite.
Internal subdivisions in the last assertion do not enter that number.
Geometric smallness is fixed before a limiting slope or any asymptotic or
derivative request is selected.
\end{theorem}

The reductions are organized by the scalar identity from
Section~\ref{sec:subdivision}.  Nozero boxes use their integrated
displacement bound.  Simple-root boxes first linearize the state:
the zero time exponent gives the two clocks \(L,W\), while a nonzero
time exponent leads to an additive action.  For annuli, the pair
\((m-1,\lambda)\) determines the initial power or logarithmic change.
The final annular case \((m,\lambda)=(1,0)\) is split by slope and
generator resonance.  In every case the ensuing clock limits select
the ordinary, one-rate, additive, or mixed transfer.

\subsection{Conventions and two elementary analytic facts}

All the analytic functions below extend to slightly larger complex boxes
than the boxes used for their arguments and trajectories.  Ordinary
parameters are suppressed from formulas; identities in profile variables
hold for all these parameters.  At each physical passage the field and
preparation parameters are held fixed while time or state is differentiated.
Their derivatives as independent arguments of a transfer are retained later.
Once an independent-argument kernel has been chosen, its field and all its
exogenous arguments are retained when a limiting regime is rewritten.
An alternative formula that agrees only after physical ties are imposed is
used as a new physical encoding, not as an identity of the old ambient
kernel.  Old ambient derivatives always use an identity on an open set of
the old arguments.
Smallness means a sufficiently small fixed geometric constant relative to
the finite exponent lists and these larger analytic boxes.  It never means a
quantity to be chosen after an expansion order.

Order the positive clock endpoints as \(t_-<t_+\) and put
\(L=\log(t_+/t_-)>0\).  A signed monomial \(b t^r\), \(r\ne0\), is exactly
\begin{equation}\label{eq:terminal-power-layer}
 b t^r=
 \begin{cases}
 (b t_-^r)e^{-|r|Ls},&r<0,\\
 (b t_+^r)e^{-|r|L(1-s)},&r>0,
 \end{cases}
 \qquad t=t_-e^{Ls},\quad 0\le s\le1.
\end{equation}
The amplitude is taken at the end where its absolute value is largest.
Zero amplitudes are allowed and are not subjected to root extraction.
Reversing the traversal interchanges the two ends.  The orientation used for
a transfer need not be positive physical time.

\begin{lemma}[Rational diagonal primitives]\label{lem:terminal-diagonal}
Let \(D_r=\sum_j r_jX_j\partial_{X_j}\), with fixed rational \(r_j\), and
let \(a(X)\) be analytic near the origin, uniformly in ordinary parameters.
For rational \(\lambda\), write \(a=\sum_\nu a_\nu X^\nu\).  Then
\begin{equation}\label{eq:terminal-diagonal}
 a=(\lambda+D_r)H+R,\qquad
 H=\sum_{\lambda+r\cdot\nu\ne0}
       \frac{a_\nu}{\lambda+r\cdot\nu}X^\nu,
 \qquad
 R=\sum_{\lambda+r\cdot\nu=0}a_\nu X^\nu
\end{equation}
is analytic on a smaller common box and \(D_rR=-\lambda R\).
For \(\lambda=0\), the decomposition is \(a=\bar a+D_rH\), with
\(D_r\bar a=0\) and \(H(0)=0\).
\end{lemma}
\begin{proof}
A common denominator for the finitely many rational numbers shows that the
nonzero denominators in \eqref{eq:terminal-diagonal} have absolute value at
least a fixed positive number.  Thus both subseries converge absolutely on
every smaller polydisk, with the same property after any fixed number of
ordinary or profile derivatives.  Termwise differentiation proves the
identities.  If \(X_j=b_jt^{r_j}\), a weight-zero monomial satisfies
\(X^\nu=b^\nu\); hence \(\bar a\) is constant for every nearby value of
\(t\), at fixed coefficients, not just along a solution of an ODE.
\end{proof}

We will also use this simple real estimate.  If a profile \(X\) has fixed
sign, \(|X|\le\eta\), and its logarithmic slope in an oriented variable
\(t\) has fixed sign and absolute value at least \(a>0\), then
\begin{equation}\label{eq:terminal-profile-integral}
 \int |X(t)|\,\dd t\le \eta/a.
\end{equation}
Indeed \(|(\dd/\dd t)|X||\ge a|X|\), and integration bounds the left side
by the total variation of \(|X|/a\).  The assertion includes the identically
zero profile by continuity.  For an analytic field vanishing when specified
profiles vanish, its norm is bounded by a constant times the sum of their
absolute values on smaller boxes.  Thus \eqref{eq:terminal-profile-integral}
controls both its integrated forcing and its integrated state Lipschitz
constant.  The integral equation, or Gronwall's inequality, then keeps a
smaller input disk in the larger state disk.  Stable linear terms improve
this estimate.

If the subdivision field is identically zero, its transfer is simply
\(z_{\rm out}=z_{\rm in}\).  This is an ordinary analytic kernel for
every clock length, with its physical endpoint coordinate graphs retained.

\subsection{Nozero boxes}

The prepared field is
\begin{equation}\label{eq:terminal-box}
 D z=B\tau^\lambda f(X,z),\qquad
 D=\frac{\dd}{\dd\log\tau},\quad
 X_j=b_j\tau^{r_j},\quad B>0,
\end{equation}
where \(r_j\ne0\) are rational, \(X\) is tiny, and constant arguments are
ordinary parameters.  In a nozero box, \(|f|\ge m_f>0\) with fixed sign.
Use small state bins of width \(h\).  Monotonicity on a passage contained in
one bin gives
\begin{equation}\label{eq:terminal-nozero-integral}
 B\int \tau^\lambda|\dd\log\tau|\le h/m_f.
\end{equation}
Choose \(h\) sufficiently small against the already fixed analytic bounds.

For \(\lambda=0\), put \(b=BL\).  The normalized physical equation is
\(z_s=b f(X,z)\), with \(b\) small by
\eqref{eq:terminal-nozero-integral}.  For a large \(L\), choose the following
independent-argument family:
\begin{equation}\label{eq:terminal-nozero-extension}
 z_s=b f(0,z)+L b'\bigl(f(X,z)-f(0,z)\bigr),
 \qquad b'=B,\quad b=Lb'\quad\hbox{on the physical graph}.
\end{equation}
The small base flow \(b f(0,z)\) is regular with room, and the perturbation
vanishes at \(X=0\) for every independent \(b,b'\).  Formula
\eqref{eq:terminal-power-layer} therefore gives the regular one-rate model
with rate \(L\).  On the long-length rule, \(b'=b/L\) is bounded and small.
If \(L\) is bounded, use \(z_s=b f(X,z)\) directly; the layer exponentials
are bounded analytic functions of \((L,s)\), including at \(L=0\).
This is the physical short rule.  At a bounded-\(L\) corner of the
already chosen independent family \eqref{eq:terminal-nozero-extension},
retain that same displayed field with independent \(b,b'\); it too is
an ordinary normalized analytic field, and hence preserves its original
ambient derivatives.

For \(\lambda\ne0\), let \(\tau_e\) be the end where \(\tau^\lambda\)
is largest and set \(b_0=B\tau_e^\lambda\).  Exactly,
\begin{equation}\label{eq:terminal-nozero-power-integral}
 B\int\tau^\lambda|\dd\log\tau|
 =\frac{b_0}{|\lambda|}\bigl(1-e^{-|\lambda|L}\bigr).
\end{equation}
Thus \(L\ge1\) makes \(b_0\) small.  The coefficient
\(B\tau^\lambda\) is itself a nonconstant endpoint layer, so the normalized
field has zero base and the form \(L W(E,H,z)\), with
\(W(0,0,z)=0\).  It is a regular one-rate transfer when \(L\) grows.
For a bounded short interval use the small bounded prefactor
\(b_1=LB\tau_e^\lambda\) and the ordinary equation
\(z_s=b_1e^{-|\lambda|L s_e}f(X,z)\), where \(s_e=s\) or \(1-s\).
All formulas extend at \(L=0\).  The short and long rules may overlap,
for example on \(1/2<L<2\); they are endpoint choices, not cuts defined by
an unbounded logarithm in the physical plane.

\subsection{Simple boxes}

Now \(f(X,z)\) in \eqref{eq:terminal-box} has a regular analytic simple
root \(z=c(X)\); the physical state stays sufficiently close to this root.
Write \(b(X)=f_z(X,c(X))\), which is nonzero with fixed sign on a root
chart.  Set
\begin{equation}\label{eq:terminal-linearizer}
 \psi(X,z)=(z-c(X))
 \exp\left(\int_{c(X)}^z
     \left\{\frac{b(X)}{f(X,\zeta)}-
                    \frac1{\zeta-c(X)}\right\}\dd\zeta\right).
\end{equation}
The apparent pole in the integrand is removable.  The formula defines an
analytic coordinate, \(\psi_z(X,c)=1\), and direct differentiation gives
\(\psi_z f=b\psi\).  With \(v=\psi(X,z)\),
\begin{equation}\label{eq:terminal-simple-linearized}
 Dv=B\tau^\lambda b(X)v+g(X,v),\qquad g(0,v)=0.
\end{equation}
Here \(g\) is \(D_X\psi\) expressed in \((X,v)\), with
\(D_X=\sum r_jX_j\partial_{X_j}\); the ordinary parameters are fixed.
Its vanishing is an analytic identity on the entire independent profile
box.  We always traverse the linear term in its damped direction.

\subsubsection{Zero time exponent}

For \(\lambda=0\), apply Lemma~\ref{lem:terminal-diagonal} to obtain
\(b=\bar b+D_XH\), \(D_X\bar b=0\), \(H(0)=0\).  The invariant
\(\bar b\) has the sign of \(b(0)\).  Define
\begin{equation}\label{eq:terminal-rho}
 \rho=\tau\exp(H/\bar b),\qquad
 Y_j=X_j\exp(r_jH/\bar b)=b_j\rho^{r_j}.
\end{equation}
The profile map has identity Jacobian at zero, hence an analytic inverse on
smaller boxes, and
\(D\log\rho=b/\bar b>0\).  In the damped oriented logarithmic variable
\(t\in[0,L]\) of \(\rho\), the exact equation is
\begin{equation}\label{eq:terminal-simple-mixed}
 v_t=-c v+\widetilde g(Y,v),\qquad
 c=B|\bar b|,\quad W=cL,
 \qquad \widetilde g(0,v)=0.
\end{equation}
The profiles are endpoint layers of fixed positive rational rates.
The normalized equation, useful also when either clock vanishes, is
\(v_s=-Wv+L\widetilde g(Y,v)\).  The complete clock alternatives are
\begin{center}
\begin{tabular}{p{0.43\linewidth}p{0.48\linewidth}}
\hline
Clock behavior along the test sequence & Exact model used \\
\hline
\(L,W\to\infty\), with \(W/L\to0\), a positive finite limit,
or \(\infty\) & Mixed model (I), with the auxiliary ratio prescribed by
Theorem~\ref{thm:mixed-packets} in the unbalanced cases. \\
\(W\to\infty\), \(L\) bounded, including \(L\to0\)
& Stable one-rate model with rate \(W\); all profiles are slow analytic
functions of \((L,s)\). \\
\(L\to\infty\), \(W\) bounded, including \(W\to0\)
& Regular one-rate model with rate \(L\), base \(-Wv\). \\
Both clocks bounded, including zero limits
& Ordinary normalized analytic transfer. \\
\hline
\end{tabular}
\end{center}
For a finite but large limiting bounded clock, subdivide the normalized
interval into a fixed number of fractions so its length per fraction is
small.  This gives bounded base-flow inverses and fixed analytic room.
Exact endpoint shifts regenerate the profiles at internal ends; damping and
\eqref{eq:terminal-profile-integral} keep the real states in the smaller
disk.  The rigorous status of this internal operation is specified below.

\subsubsection{Nonzero time exponent and the additive action}

For \(\lambda\ne0\), diagonal inversion gives
\begin{equation}\label{eq:terminal-simple-diagonal}
 b=(\lambda+D_X)H+R,\qquad D_XR=-\lambda R,
 \qquad H_0=H(0)=b(0)/\lambda\ne0,
 \quad R(0)=0.
\end{equation}
Solve
\begin{equation}\label{eq:terminal-h}
 H=H_0h^\lambda+R\log h
\end{equation}
for \(h\) near 1.  The derivative in \(h\) at the origin is
\(\lambda H_0=b(0)\ne0\), so this is a bounded analytic operation.
Put
\begin{equation}\label{eq:terminal-simple-w}
 \begin{gathered}
 I=B\tau^\lambda R,\quad
 w=|BH_0|(\tau h)^\lambda,
 \quad\sigma=\sgn(BH_0),\quad
 \gamma=\sigma I/\lambda,\\
 Y_j=X_jh^{r_j},\qquad
 \frac{\gamma}{w}=\frac{R}{b(0)h^\lambda}.
 \end{gathered}
\end{equation}
Here \(DI=0\), \(Y_j\) is a constant multiple of \(w^{r_j/\lambda}\),
and \(X\mapsto Y\) is an analytic change tangent to the identity.
The exact primitive identity is
\begin{align}\label{eq:terminal-exact-primitive}
 B\tau^\lambda H+I\log\tau
 &=BH_0(\tau h)^\lambda+I\log(\tau h)\notag\\
 &=\sigma(w+\gamma\log w)-\frac I\lambda\log|BH_0|.
\end{align}
The last term is constant in time.  Also
\begin{equation}\label{eq:terminal-clock-jacobian}
 D\log w=\frac{\lambda b(X)}{b(0)h^\lambda+R},
\end{equation}
a nonzero analytic factor near \(\lambda\).  Differentiating
\eqref{eq:terminal-exact-primitive} therefore turns
\eqref{eq:terminal-simple-linearized} into
\begin{equation}\label{eq:terminal-additive-a}
 v_w=\sigma(1+\gamma/w)v+w^{-1}\widehat g(Y,v),
 \qquad \widehat g(0,v)=0.
\end{equation}
Thus no sign or factor of \(\lambda\) is absorbed in an asymptotic
replacement.

On the large-action part, choose \(q_*\) sufficiently large once for these
finite formulas and use
\begin{equation}\label{eq:terminal-action}
 \begin{gathered}
 q=\min w\ge q_*,\quad p=\max w,\quad
 \delta=\gamma/q,\\
 A(w)=w+\gamma\log(w/q),\quad
 Q=A(p),\quad S=Q-q>0.
 \end{gathered}
\end{equation}
The tiny bound on \(\gamma/w\) makes \(A'\) positive and bounded away
from zero.  In particular \(p\asymp Q\) and
\(Q/p=1+O(\delta)\).  Negative profile powers are anchored at \(q\);
positive powers may be anchored at \(Q\), with their exact factor
\((w/Q)^r\).  Their amplitudes are still tiny because \(Q/p\) is close
to 1.  Include \(\gamma/w\) as a negative profile.  Then
\eqref{eq:terminal-additive-a} is exactly additive model (A), in its damped
direction.

If \(Q/q\to\infty\), apply Theorem~\ref{thm:additive-packets}.
For the complementary bounded-ratio case set
\begin{equation}\label{eq:terminal-bounded-action}
 M=S/q,\quad s=(A-q)/S,\quad y=w/q,
 \qquad y+\delta\log y=1+Ms.
\end{equation}
Its implicit derivative is \(1+\delta/y\), a unit on the real interval.
The normalized forward equation is exactly
\begin{equation}\label{eq:terminal-bounded-a}
 v_s=\sigma S v+\frac{M}{y+\delta}\widehat g(Y,v).
\end{equation}
Reverse \(s\) if required for damping.  The profiles are analytic powers
of \(y\) and \(y/(1+M)\).  For bounded \(M\), this is a stable one-rate
model if \(S\to\infty\) and an ordinary transfer if \(S\) is bounded,
including \(M=0\) or \(S=0\) in its normalized extension.  A large finite
limit of \(M\) is treated by the fractions in
\eqref{eq:terminal-local-ratio} below.

For the remaining bounded range of \(w\), \(B\tau^\lambda\) is bounded
because \(h\) and \(H_0\) have fixed unit bounds.  Cut using a tiny
positive threshold for the subanalytic quantity \(B\tau^\lambda\).
Above it, the logarithmic length is uniformly bounded: both endpoint values
of \(B\tau^\lambda\) lie between two fixed positive constants, and their
log ratio equals \(\lambda\) times the clock length.  This gives an
ordinary transfer.  Below it, \(B\tau^\lambda\) is an additional tiny
nonconstant layer.  Both terms of
\eqref{eq:terminal-simple-linearized} then vanish without layers; use a
regular one-rate transfer for large log length and the ordinary normalized
transfer otherwise.  These cuts involve no unbounded logarithm.

\subsection{Annuli with at least one nonzero power exponent}

The annular identity is
\begin{equation}\label{eq:terminal-annulus}
 Du=A_0\tau^\lambda u^m F(X),\qquad
 X_j=\chi_j\tau^{a_j}u^{\beta_j},\qquad u>0,
\end{equation}
where \(A_0\ne0\) is parameter-only, \(m\) is integral, all exponents
are rational, and \((a_j,\beta_j)\ne(0,0)\).
The profiles are tiny and \(F(0)=1\), so \(F>0\) and
\(\partial_{\log u}F\) is small.  Pure multipliers in the following
changes remain parameter-only until an endpoint normalization is explicitly
made.

\subsubsection{Two nonzero power exponents}

Suppose \(m\ne1\) and \(\lambda\ne0\).  Set
\begin{equation}\label{eq:terminal-power-swap}
 U=u^{1-m},\qquad T=\tau^\lambda,\qquad
 \frac{\dd U}{\dd T}=kF,
 \qquad k=\frac{(1-m)A_0}{\lambda}.
\end{equation}
Split by fixed tiny and huge thresholds of \(|k|T/U\).
In the tiny range choose \(U_*\) at \(T_{\max}\).  Since \(F\) is
bounded,
\[
 |U(T)-U_*|\le C|k|T_{\max},\qquad
 |k|T_{\max}/U_*\ll1.
\]
Thus \(V=U/U_*\) is close to 1.  Its field in \(\log T\) is
\[
 V_{\log T}=(kT/U_*)F.
\]
The first factor is a tiny endpoint layer.  Every original profile is an
end-normalized rational power of \(T\), times an analytic power of
\(V\); a zero time power is a tiny ordinary argument.  Use the zero-base
regular one-rate model or the bounded-length ordinary model.  The defining
conjugacy holds for every nearby positive choice of \(U_*\).

In the huge range invert \eqref{eq:terminal-power-swap}:
\(\dd T/\dd U=k^{-1}F^{-1}\).  The clock \(U\) is transverse since
\(kF\ne0\), and the new ratio \(|k|^{-1}U/T\) is tiny, reducing to the
previous case.  In the comparable range put \(z=U/(|k|T)\).  It stays in
a fixed compact positive interval, and
\begin{equation}\label{eq:terminal-comparable-power}
 z_{\log T}=\sgn(k)F-z.
\end{equation}
The profiles become pure time powers times analytic state powers.  The
state derivative of the right side is \(-1+o(1)\), uniformly on this
compact interval after the already permitted generator smallness choice.
Choose a fixed small threshold for the right side.  Away from zero this is
a nozero box; near zero the nonzero derivative gives a unique simple-root
chart with room.  Both were treated above.

\subsubsection{Exactly one zero power exponent}

If \(m=1,\lambda\ne0\), use
\[
 v=\log u,\quad w_0=|A_0/\lambda|\tau^\lambda,
 \quad \kappa=\sgn(A_0/\lambda).
\]
If \(m\ne1,\lambda=0\), invert the transverse field and use
\[
 v=\log\tau,\quad
 w_0=u^{1-m}/|(1-m)A_0|,
 \quad\kappa=\sgn((1-m)A_0),
\]
replacing \(F\) by \(F^{-1}\).  In either case the exact equation is
\begin{equation}\label{eq:terminal-log-power}
 v_{w_0}=\kappa F(X),\qquad
 X_i=\widetilde\chi_iw_0^{p_i}e^{d_i v},\quad\kappa\in\{-1,1\},
\end{equation}
with fixed rational \((p_i,d_i)\ne(0,0)\).
Absolute \(v\) is not an argument of a restricted analytic function;
the underlying endpoint variable is the positive \(e^v\).

For bounded \(w_0\), take finitely many small absolute-width bins.
Then \(v-v_-\) is small, and in \(\log w_0\) its derivative is
\(\kappa w_0F\).  Write each profile using its large time-power end,
times \(e^{d_i(v-v_-)}\).  The amplitudes are tiny, since the state
displacement is small and the original endpoint profiles are tiny.
Near \(w_0=0\), the coefficient \(w_0\) is an extra tiny layer; away
from zero the logarithmic length is bounded.  These give respectively the
regular one-rate and ordinary rules, with the bounded-length overlap.
For the ordinary rule keep an explicit integrated-coefficient guard in
the independent family.  Let \(w_e=w_{0,+}\) and
\(L=\log(w_{0,+}/w_{0,-})\).  The exact normalized field is
\begin{equation}\label{eq:terminal-bounded-w-guard}
 \begin{gathered}
 z_s=\kappa Lw_e e^{-L(1-s)}F,\qquad z=v-v_-,\\
 \int_0^1 Lw_e e^{-L(1-s)}\,\dd s
 =w_e(1-e^{-L})=w_{0,+}-w_{0,-}.
 \end{gathered}
\end{equation}
Impose \(w_e(1-e^{-L})<h\), with \(h\) chosen so that the forcing
and state Lipschitz bounds \(Ch\), even on the closed range, keep
solutions strictly inside the larger state disk.  This is an analytic
inequality in the old independent \((L,w_e)\) and their affine
exponential, so it preserves the original kernel and its ambient
derivatives.  The physical absolute-width bins satisfy this guard.
All other profiles retain their independent endpoint amplitudes and
the analytic state factor \(e^{d_i(v-v_-)}\).  Bounded length alone
would not supply this estimate.  Near zero the tiny coefficient-layer
bound already controls the integral by \(w_e\), independently of the
logarithmic length.

For large \(w_0\), separate the pure profiles \(X^0\), those with
\(d_i=0\), and set \(F^0=F(X^0,0)\).  Apply
Lemma~\ref{lem:terminal-diagonal} with weight 1:
\begin{equation}\label{eq:terminal-annular-diagonal}
 F^0=(1+D_p)H+R,\quad D_pR=-R,\quad H(0)=1,
 \qquad H=h+R\log h,
 \quad D_p=\sum_{d_i=0}p_iX_i\partial_{X_i}.
\end{equation}
Set \(w=w_0h\), \(\gamma=w_0R\), and \(Y_i=X_i^0h^{p_i}\).
The profile change is regular, \(\gamma\) is constant on the time line,
and
\begin{equation}\label{eq:terminal-annular-primitive}
 \int F^0\,\dd w_0=w+\gamma\log w+\text{constant},\qquad
 \frac{\dd w}{\dd w_0}=\frac{F^0}{1+\gamma/w}.
\end{equation}
Use \(q,p,\delta,A,Q,S\) from \eqref{eq:terminal-action}, choosing the
cut high enough that \(q\ge q_*\).  From the smaller-action end put
\begin{equation}\label{eq:terminal-annular-residual}
 z=v-v_- -\kappa(A(w)-q),\qquad
 z_w=\kappa\frac{\dd w_0}{\dd w}\bigl(F-F^0\bigr),\quad z(q)=0.
\end{equation}
For a nonpure nonzero profile, its logarithmic slope in \(w_0\) is
\[
 \frac{p_i}{w_0}+d_i\kappa F.
\]
It has the sign of \(d_i\kappa\) and absolute value bounded below by a
fixed multiple of \(|d_i|\) once the cutoff is large and the generators
are small.  Equation~\eqref{eq:terminal-profile-integral} therefore bounds
the total residual displacement by a fixed constant times the amplitudes.

For \(d_i\kappa<0\) and \(d_i\kappa>0\), respectively, define
\begin{align}\label{eq:terminal-annular-amplitudes}
 e_i&=\widetilde\chi_iq^{p_i}e^{d_iv_-},&
 E_i&=e_i(w/q)^{p_i}e^{-|d_i|(A-q)},\notag\\
 D_i&=\widetilde\chi_iQ^{p_i}e^{d_iv_+-d_iz_+},&
 H_i&=D_i(w/Q)^{p_i}e^{-|d_i|(Q-A)}.
\end{align}
In the second line tie the proposed output residual by
\begin{equation}\label{eq:terminal-additive-end-tie}
 e^{v_+}=e^{v_-}\exp(\kappa S+z_+).
\end{equation}
Then the original profile is exactly \(E_i h(Y)^{-p_i}e^{d_i z}\) or
\(H_i h(Y)^{-p_i}e^{d_i z}\).  In particular the residual field is
analytic in pure profiles, the fast profiles, and \(z\), and vanishes
when all fast profiles vanish.  The rates \(|d_i|\) belong to one positive
rational lattice.  This is additive model (B), with tiny right amplitudes
because \(Q/p\) and \(h\) are bounded units and \(z_+\) is tiny near
the physical passage.

For \(Q/q\to\infty\), use its additive packets.  For bounded
\(M=S/q\), the action variable in \eqref{eq:terminal-bounded-action}
gives instead
\begin{equation}\label{eq:terminal-bounded-b}
 z_s=S\frac{y}{y+\delta}G(X,E,H,z),
 \qquad E_i=\widehat e_i(s)e^{-|d_i|Ss},\quad
 H_i=\widehat D_i(s)e^{-|d_i|S(1-s)},
\end{equation}
where the hatted amplitudes are the original amplitudes times analytic
slow powers of \(y\) and \(y/(1+M)\).  This is the zero-base regular
one-rate model if \(S\to\infty\), or an ordinary transfer if \(S\)
is bounded.  Both bounded-ratio models, (A) and (B), allow \(M\to0\).

Here are explicit fractions when the finite limit of \(M\) is large.
Write \(A=q+S(j+t)/K\), \(0\le t\le1\), and let \(y_0\) be the value
of \(y\) at \(j/K\).  Put \(y=y_0R\).  Subtracting the equation for
\(y_0\) from \eqref{eq:terminal-bounded-action} gives
\begin{equation}\label{eq:terminal-local-ratio}
 R+\frac{\delta}{y_0}\log R
 =1+\frac{M}{Ky_0}t,\qquad y_0\ge1.
\end{equation}
Choose one finite \(K\) making \(M/K\) small on a neighborhood of the
chosen finite limit.  Then \(R\) and every fixed power of \(R\) have
fixed analytic room near 1.  An internal left amplitude is, exactly,
\[
 e_i^{(j)}=e_i(w_j/q)^{p_i}e^{-|d_i|jS/K},
\]
and there is the analogous right formula.  Adverse powers are paid by the
exponential: the logarithmic derivative in \(A\) of the absolute left
factor is
\[
 \frac{p_i}{w+\gamma}-|d_i|
 \le\frac{|p_i|}{q_*(1-|\delta|)}-|d_i|<0
\]
after increasing the fixed \(q_*\).  The reversed estimate handles the
right factor.  Thus internal amplitudes remain tiny, and pure powers are
anchored at their local large ends.  The total real forcing estimate is
unchanged.  Every new argument uses bounded regular solves, fixed rational
powers, and exponentials of fixed fractions of \(S\).

\subsection{The logarithmic annulus and its small-slope boundary}

It remains to treat \(m=1,\lambda=0\).  Set
\begin{equation}\label{eq:terminal-log-annulus}
 \theta=\log\tau,\quad v=\log u,\qquad
 v_\theta=A_0F(X),\quad X_i=\chi_i e^{a_i\theta+\beta_i v}.
\end{equation}
For huge \(|A_0|\), exchange the two transverse logarithmic variables:
\(\dd\theta/\dd v=A_0^{-1}F^{-1}\).  This interchanges the rational
weights and enters the small-coefficient case below.  Hence it suffices to
consider bounded \(|A_0|\).

\subsubsection{Slopes separated from zero and resonance}

Exclude a fixed neighborhood of zero and of every nonzero number
\(\alpha=-a_i/\beta_i\) with \(\beta_i\ne0\).  On each resulting
coefficient interval, \(|a_i+\beta_iA_0|\) has a fixed positive lower
bound.  Since \(F-1\) is tiny, the physical logarithmic slopes
\(a_i+\beta_iA_0F\) have the same signs and fixed lower bounds.  It
follows that all profiles have small integral in \(\theta\).
Put \(\kappa=\sgn A_0\), \(c=|A_0|\),
\(L=\theta_+-\theta_->0\), \(W=cL\), and
\[
 z(t)=v(\theta_-+t)-v_- -\kappa ct,\qquad 0\le t\le L.
\]
Its field is \(z_t=A_0(F-1)\), and \(z\) stays small.  Write
\(\rho_i=a_i+\beta_i\kappa c\), with its fixed nonzero sign.  For
\(\rho_i<0\) use a left amplitude; for \(\rho_i>0\) use a right
amplitude.  The resulting positive rate is \(|\rho_i|\), an affine
rational function of the exact occurrence \(c=W/L\).  The factor in
state is \(e^{\beta_i z}\), and the endpoint residual tie is
\begin{equation}\label{eq:terminal-mixed-end-tie}
 e^{v_+}=e^{v_-}\exp(\kappa W+z_+).
\end{equation}
The same cancellation as in \eqref{eq:terminal-annular-amplitudes}
gives mixed model (II).  Bounded occurrences of \(A_0\) in its field
may be independent ordinary parameters, tied separately.  If \(L\)
grows, then \(W\) grows and \(c\) stays in a compact positive interval;
if \(L\) is bounded, so is \(W\), and the normalized transfer is
ordinary, including zero length.
In this independent kernel domain impose the selected bounded slope
interval on the exact \(W/L\), with slightly relaxed endpoints still
separated from zero and the resonances.  These are linear inequalities
between \(W\) and \(L\), in addition to the fixed positive rate margins.
An independent ordinary copy of \(A_0\) does not replace these
inequalities.

\subsubsection{Straightening the pure-state field}

For sufficiently small \(|A_0|>0\), let \(X^s\) denote the profiles
with \(a_i=0\) and set \(F_s=F(X^s,0)\).  The pure-state weights
\(\beta_i\) here are nonzero.  Apply diagonal inversion to the reciprocal:
\begin{equation}\label{eq:terminal-state-diagonal}
 F_s^{-1}=b+D_\beta H,\quad D_\beta b=0,\quad
 b(0)=1,\quad H(0)=0,\qquad
 D_\beta=\sum_{a_i=0}\beta_iX_i\partial_{X_i}.
\end{equation}
The following invariance is needed for mismatched endpoint values.  If
\(\beta\cdot\nu=0\), then for every nearby \(v\), at fixed
coefficients,
\begin{equation}\label{eq:terminal-state-invariant}
 (X^s(v))^\nu=\chi^\nu e^{(\beta\cdot\nu)v}=\chi^\nu.
\end{equation}
Consequently \(b(X^s(v))\) is independent of both \(v\) and \(\theta\)
on the physical profile surface.  One can compute it at either proposed
endpoint without imposing a redundant equality between two evaluations.
This does not make its derivatives in independent coefficients or profiles
zero.

Define
\begin{equation}\label{eq:terminal-state-change}
 \widetilde v=v+H/b,\qquad
 Y_i=X_i\exp(\beta_iH/b)
      =\chi_i e^{a_i\theta+\beta_i\widetilde v}.
\end{equation}
The pure-state part of the profile map has identity Jacobian at zero;
invert it analytically first and then invert the remaining profiles by
their unit factors.  Direct differentiation, using
\eqref{eq:terminal-state-diagonal}, gives
\begin{equation}\label{eq:terminal-state-jacobian}
 \widetilde v_v=1+\frac{D_\beta H}{b}=\frac1{bF_s},
 \qquad
 \frac{\dd\widetilde v}{\dd\theta}
 =\frac{A_0}{b}\frac{F}{F_s}.
\end{equation}
Thus this is an actual state diffeomorphism.  In the positive coordinate
\(R=e^{\widetilde v}=u\exp(H/b)\), its derivative is
\[
 R_u=\frac{\exp(H/b)}{bF_s}>0.
\]
Let \(p=A_0/b=\kappa c\), where \(c>0\), and subtract the exact slope:
\begin{equation}\label{eq:terminal-small-residual}
 z=\widetilde v-\widetilde v_- -\kappa ct,
 \qquad z_t=p\left(\frac{F}{F_s}-1\right).
\end{equation}
Under the analytic inverse profile change, \(F/F_s-1\) vanishes
identically when all profiles with \(a_i\ne0\) vanish.  Hence every
nonzero residual Taylor monomial contains such a profile.

Here the small-slope threshold can be fixed from the finite weight list.
For example, when there are nonzero \(a_i\), take
\[
 c_*<\frac{\min_{a_i\ne0}|a_i|}
                 {4(1+\max_i|\beta_i|)}.
\]
By first bounding \(b,F,F_s\) near 1 and choosing \(|A_0|\) small enough,
both the straight-line rates and the physical logarithmic slopes of these
profiles have the sign of \(a_i\) and magnitude at least
\(|a_i|/2\).  If no such profile exists, the residual is identically zero
and no threshold is needed.  Pure-state layers have positive rates
\(|\beta_i|c\), with their end selected by the sign of
\(\beta_i\kappa\).  The residual has small total integral since every
monomial contains a uniformly fast profile; the other profiles stay tiny.
More explicitly, on a smaller fixed box the analytic vanishing gives
\[
 |G|+|G_z|\le C\sum_{a_i\ne0}|Y_i|,
 \qquad
 \int_0^L(|G|+|G_z|)\,\dd t
 \le \frac{CN\eta}{a_*},
\]
where \(G=p(F/F_s-1)\), the \(N\) designated fast profiles have amplitudes
at most \(\eta\), and their rates are at least a fixed \(a_*>0\).
The bounds are uniform in the remaining tiny profiles and the bounded
coefficient \(p\).  Thus \(\eta\) can be fixed using only analytic room
and the finite fast-rate margin.  No integral of an isolated pure-state
layer of rate \(|\beta_i|c\) is used in this bound.
For \(L,W=cL\to\infty\), this is mixed model (II), including
\(c\to0\).  A positive but arbitrarily small limiting \(c\) requires
no new geometric smallness: the designated fast factor has a uniform
positive rate, exactly the uniformity clause of
Theorem~\ref{thm:mixed-packets}.
The independent domain imposes \(0<W<c_*L\) itself, as well as the
bounded range for the independent coefficient \(p\).  Thus the exact
slope stays in the selected small-slope regime even away from the
physical tie \(pL=\kappa W\).

For completeness, when \(L\to\infty\) and \(W\) is bounded, put
\(t=Ls\).  For \(a_i<0\) the left layer factors exactly as
\begin{equation}\label{eq:terminal-signed-slow-left}
 e^{(a_i+\beta_i\kappa c)Ls}
 =e^{-|a_i|Ls}\,e^{\beta_i\kappa Ws},
\end{equation}
whereas for \(a_i>0\) the right layer factors as
\begin{equation}\label{eq:terminal-signed-slow-right}
 e^{-(a_i+\beta_i\kappa c)L(1-s)}
 =e^{-|a_i|L(1-s)}\,e^{-\beta_i\kappa W(1-s)}.
\end{equation}
These slow factors may grow or decay: their signs are retained.  Pure-state
layers are \(e^{-|\beta_i|Ws}\) or
\(e^{-|\beta_i|W(1-s)}\), also slow.  They are ordinary bounded
analytic arguments at a finite \(W\) limit.  The field
\(Lp(F/F_s-1)\) has a uniformly fast layer in each monomial, so it is
a zero-base regular one-rate field of rate \(L\), with \(p\) an
independent bounded coefficient.  If necessary, choose fixed fractions
making the bounded \(W\)-length per fraction small; regenerate amplitudes
with the full affine shift \(a_iL+\beta_i\kappa W\).  This avoids
enlarging a tiny fast amplitude by a large finite slow exponential.
If \(L\) is bounded, \(W\) is bounded as well, and the normalized
formula is ordinary, including \((L,W)=(0,0)\).  Thus the small-slope
alternatives are exhaustive:
\begin{center}
\begin{tabular}{p{0.43\linewidth}p{0.48\linewidth}}
\hline
Clock behavior & Model \\
\hline
\(L,W\to\infty\), \(0\le\lim W/L\le c_*\)
& Mixed (II); at zero slope every residual monomial has a fast factor.\\
\(L\to\infty\), bounded \(W\), including \(W\to0\)
& Regular one-rate, using the signed factorizations above.\\
Bounded \(L\) (hence bounded \(W\)), including zero limits
& Ordinary normalized transfer.\\
\hline
\end{tabular}
\end{center}
The residual tie is
\(R_+=R_-\exp(\kappa W+z_+)\).  It has the same right-amplitude
cancellation as \eqref{eq:terminal-mixed-end-tie}.

\subsubsection{Nonzero resonances}

For each nonzero rational generator resonance
\(\alpha=-a_i/\beta_i\), weight the physical state before solving the
passage:
\begin{equation}\label{eq:terminal-resonance-weight}
 U=u/\tau^\alpha,\qquad DU=U(A_0F-\alpha).
\end{equation}
This is still an exact rational field in \(U\) with subanalytic base
coefficients.  Its state derivative is
\begin{equation}\label{eq:terminal-resonance-derivative}
 \partial_U(DU)=A_0F-\alpha+A_0\partial_{\log u}F.
\end{equation}
It is as small as prescribed by the resonance width and generator
smallness.  Apply the initial rational subdivision again.
Lemma~\ref{lem:subdivision-resonance-redo} proves that this redo terminates:
the finite candidate exponent lists are prepared before choosing those
widths; only a new initial unscaled annulus with powers \((m',\lambda')
=(1,0)\) could encounter this issue again, and its leading slope is
bounded by a fixed multiple of
\eqref{eq:terminal-resonance-derivative}.  Choose the original widths to
place it below the small-slope threshold for every precomputed list.
The new formula's own generator cuts then put its units near 1.  Contact
bins have \(\lambda'=1\), and outer annuli in the box induction have
\(m'\ge2\), so neither can restart a nonzero-resonance redo.  The huge
slope swap already enters the small case and is not used in this redo.

\subsection{Endpoint ties, extensions, and differentiated identities}

\begin{lemma}[Identity on the nearby tied endpoint graph]
\label{lem:terminal-tied-endpoints}
Every transfer constructed above computes the actual scalar mismatch for
all nearby proposed endpoints on its argument graph.  The terminal state
coordinate has nonzero transverse derivative.  Endpoint-dependent
normalizations and the pure-state invariant do not destroy this assertion.
\end{lemma}
\begin{proof}
All time and profile primitives above are identities of analytic functions
on the prepared physical piece.  Invariants in a time change are constant
for every nearby time, by the weight-zero identity in
Lemma~\ref{lem:terminal-diagonal}.  The pure-state invariant is constant
for every nearby state by \eqref{eq:terminal-state-invariant}.  The
linearizer, power maps, and profile inversions have nonzero Jacobian on
their specified domains.  Clock swaps use a nonzero scalar velocity.

The only apparent dependence of the computed flow on a proposed final
state occurs in right amplitudes.  For the additive annulus, the tie
\eqref{eq:terminal-additive-end-tie} is equivalent, on its positive real
domain, to
\(z_+=v_+-v_- -\kappa S\).  If \(d_i\kappa=|d_i|\), it gives
\begin{equation}\label{eq:terminal-cancel-additive}
 D_i=\widetilde\chi_iQ^{p_i}
       \exp(d_iv_-+|d_i|S).
\end{equation}
This expression is independent of the proposed \(v_+\).  Substitution
in \eqref{eq:terminal-annular-amplitudes}, at every intermediate state
\(z\), gives
\[
 H_i h^{-p_i}e^{d_i z}
 =\widetilde\chi_iw_0^{p_i}
       \exp\bigl(d_i[v_-+\kappa(A-q)+z]\bigr),
\]
exactly the original profile.  The left identity is immediate.  Using
\(Q\), rather than \(p\), in the amplitude is exact because the layer
contains \((w/Q)^{p_i}\).

For a mixed right layer with positive
\(\rho_i=a_i+\beta_i\kappa c\), use
\[
 D_i=\chi_i\exp(a_i\theta_++\beta_iv_+-\beta_i z_+).
\]
The tie \(z_+=v_+-v_- -\kappa W\) gives
\[
 D_i=\chi_i\exp(a_i\theta_++\beta_iv_-+\beta_i\kappa W).
\]
Multiplication by \(e^{-\rho_i(L-t)}e^{\beta_i z}\) recovers
\(\chi_i e^{a_i(\theta_-+t)+\beta_i(v_-+\kappa ct+z)}\).
The calculation uses no sign restriction on \(\beta_i\) or \(\kappa\).
For the small-slope case replace \(v\) by \(\widetilde v\), which is
regular by \eqref{eq:terminal-state-jacobian}; invariance makes its
\(p=A_0/b\) independent of the proposed other endpoint.  Thus the flow
side is independent of that proposed endpoint before it is equated with
\(z_+\).

For clarity, an endpoint normalization such as \(U_*\) obeys a slightly
different but equally exact rule.  Let \(C_a\) be its regular coordinate
map, with \(a\) held temporarily independent.  The transfer is a
conjugacy for every nearby fixed \(a\), so its mismatch is
\(C_a(U_{\rm flow})-C_a(U_+)\).  At equality,
\[
 \dd\{C_a(U_{\rm flow})-C_a(U_+)\}
 =C_a'(U_+)\,\dd(U_{\rm flow}-U_+).
\]
The two \(a\)-variations cancel even when \(a\) is then chosen from a
proposed endpoint.  For a clock swap, the nonzero clock velocity likewise
determines the hitting-time variation; the remaining state derivative is
transverse and nonzero.  This proves the assertion.
\end{proof}

\begin{example}[An exact check of the right-amplitude correction]
For \(v_w=1+\chi e^v\), with no pure profiles and \(S=p-q\), the
residual equation is \(z'=D e^{-(p-w)}e^z\), where
\(D=\chi e^{v_+-z_+}=\chi e^{v_-+S}\).  Starting with \(z(q)=0\),
\[
 z(p)=-\log\bigl(1-\chi e^{v_-}(e^S-1)\bigr).
\]
Hence the derivative of the mismatch \(z(p)-z_+\) in the proposed
\(v_+\) is exactly \(-1\).  The bounded correction
\(e^{-z_+}\) in the right amplitude is necessary for this identity.
\end{example}

\begin{lemma}[Structural extensions and the physical chain rule]
\label{lem:terminal-ambient-extension}
The terminal kernels can be chosen on independent-argument neighborhoods
so that every layer-vanishing identity required by its primitive theorem
holds throughout that neighborhood.  On their physical graph they give
the stated transfer identity.  Every fixed finite derivative of the
physical identity, or of an old kernel rewritten in a new regime, is a
finite expression in derivatives of these kernels, bounded analytic
operations, fixed affine exponential shifts, and algebraic factors with
nonzero divisors on the domain.
\end{lemma}
\begin{proof}
Use the displayed structural formulas, not arbitrary extensions of a
quantity which vanishes only on the physical graph.  For nozero boxes the
choice is \eqref{eq:terminal-nozero-extension}.  If its transfer is
\(K(L,b,b',\eta)\), where \(\eta\) denotes independent amplitudes and
other arguments, the physical identity at fixed \(\eta\) is
\[
 T(L,B,\eta)=K(L,BL,B,\eta),\qquad
 T_B=LK_b+K_{b'},\quad T_L=K_L+BK_b.
\]
When the amplitudes also vary, their chain terms must be added.
For simple boxes \(g(0,v)=0\) is an analytic identity in the root-chart
parameters.  Use damping exactly \(-W/L\) in the mixed kernel and tie
\(W=B|\bar b|L\); duplicate bounded coefficients remain independent.
The other changes preserve the vanishing identically by analytic
substitution and subtraction of their pure-profile restrictions.

In particular, the admissible independent-argument small-slope residual is
\begin{equation}\label{eq:terminal-structural-extension}
 z_t=p\left(F/F_s-1\right),
\end{equation}
with \(p\) an independent bounded parameter.  Layer rates use the exact
\(c=W/L\), and the physical graph includes \(p=\kappa W/L\).
Its fast-layer vanishing holds for every independent \(p\).
The expression \(pF/F_s-\kappa W/L\) is a different ambient extension;
it has an unwanted constant term away from the graph and is not used.
For example, when \(F=F_s=1\), the chosen residual is identically zero,
whereas that other expression integrates to \(pL-\kappa W\).
Physical reconstruction adds \(\kappa W\); on the tie
\(W=\kappa pL\), its \(p\)-derivative is the correct \(L\).
Thus equality on a graph does not assert equality of unrelated ambient
partial derivatives.

More generally, suppose a rewriting is the exact identity
\[
 \Phi(a)=\Psi(a,G(a)),\qquad E(a,G(a))=0,
 \qquad\det E_y\ne0
\]
on an open neighborhood of the old arguments.  Then
\[
 G_{a_j}=-E_y^{-1}E_{a_j},\qquad
 \Phi_{a_j}=\Psi_{a_j}+\Psi_yG_{a_j}.
\]
Repeated differentiation uses finitely many derivatives and finite powers
of \((\det E_y)^{-1}\).  These denominators may be cleared while
retaining their nonzero domain conditions.  For a concrete instance, the
bounded-action identity, with all original amplitudes kept independent, is
\[
 \Phi(Q,q,\delta,\eta)
 =\Psi\bigl(Q-q,(Q-q)/q,\delta,\eta\bigr).
\]
It implies, holding \(\delta,\eta\) fixed,
\begin{equation}\label{eq:terminal-additive-chain-rule}
 \Phi_Q=\Psi_S+q^{-1}\Psi_M,
 \qquad
 \Phi_q=-\Psi_S-(Q/q^2)\Psi_M.
\end{equation}
In particular the ratio derivative terms cannot be suppressed.
The implicit function in \eqref{eq:terminal-bounded-action} has a unit
denominator and so obeys the same rule.

Positive rational powers are tied by positive roots of polynomial
equations after clearing integer powers.  These are regular on the
positive domain; signed multipliers are used algebraically.  Clock logs
are tied by
\begin{equation}\label{eq:terminal-log-ties}
 t_-=t_+e^{-L},\quad L>0;
 \qquad q=pe^{-L_w},\quad Q=p+\delta qL_w.
\end{equation}
Log-state residuals use the positive endpoint ties
\eqref{eq:terminal-additive-end-tie} and
\eqref{eq:terminal-mixed-end-tie}, placing the decaying exponential on
the appropriate side when necessary.  Thus no absolute unbounded state
logarithm is an analytic-box argument.  Mixed layer exponents are the
fixed affine combinations \(\alpha L+\beta W\); additive fractions
use fixed multiples of \(S=Q-q\).  Unbounded preparation constants are
only algebraic multipliers or graph coordinates.  They are not inserted
into a bounded analytic argument without normalization.

Internal subdivisions are identities of actual flows on an open old
argument neighborhood.  The new amplitudes use precisely these powers and
affine shifts.  Their intermediate states are fixed successively by
equations \(z_j-\Phi_j(a_j,z_{j-1})=0\); the state Jacobian is triangular
with diagonal 1.  Differentiating the exact composition identity supplies
the corresponding finite jet graphs with the same property; see
Lemma~\ref{lem:counting-jet-lift}.  Therefore their addition preserves
isolation and gives the claimed derivative formulas.  This uses the
differentiable packet theorem for each output, not an assertion that its
variational equation has the original layer-vanishing property.
\end{proof}

\begin{lemma}[Boundary corners of the chosen domains]
\label{lem:terminal-boundary-corners}
The independent-argument domains in this section can be chosen so that the
following holds for every sequence in any one of them, including a sequence
on which imposed strict margins vanish.  After passage to a subsequence
and the finite graph additions in the displayed regime reductions, either
there is a divergent enlarged coordinate to which the packet alternative
applies, or the normalized kernel factors extend analytically near all
their limiting arguments and retain state room.  This applies to the old
independent kernel and all its fixed finite derivatives through the exact
identities of Lemma~\ref{lem:terminal-ambient-extension}.  Auxiliary
algebraic graphs and equations obtained by clearing their nonzero divisors
are analytic at bounded limits; their regularity is required at the
sequence points, not at the limiting boundary point.
\end{lemma}
\begin{proof}
There are two different kinds of margin.  A \emph{working-box margin}
restricts data to a smaller part of a region on which the analytic formula
and its existence estimates are already valid.  A \emph{structural margin}
chooses a positive coordinate, a nonzero algebraic divisor, an ordering,
or a clock regime.  We check their boundaries separately.  This is a
property of the concrete domains, stronger than completeness for sequences
whose margins stay positive.

\paragraph{Analytic boxes, amplitudes, and state room.}
Choose three nested closed boxes for each bounded argument: an admitted
box, a larger working box, and an analytic-extension box, each contained in
the interior of the next.  The finite preparation and root charts provide
the outer boxes before the admitted ones are chosen.  Similarly choose
an admitted input-state set strictly inside a working state disk and the
working disk strictly inside the field's analytic state disk.  Choose the
admitted amplitude bound \(\eta\), coefficient bounds, and any integrated
prefactor bounds so that the real existence estimate holds on a slightly
larger closed range than the admitted one.  These choices are made once.
For example a zero-base passage with integrated forcing and Lipschitz
bounds \(C\eta\) sends \(|z_{\rm in}|\le r_0\) into
\[
 |z(t)|\le (r_0+C\eta)e^{C\eta}\le r_1<r_2,
\]
where \(r_2\) is the analytic state radius.  Stable damping permits the
same bound; the regular base is first removed on its specified larger
flow box.  Thus a working-box margin tending to zero reaches the edge of
the admitted box, which still has uniform analytic room.  It does not
mean that a trajectory reaches the singular boundary of the field.

Here all needed integrated estimates hold for the independent ambient
families.  In \eqref{eq:terminal-nozero-extension}, the base has the small
coefficient \(b\), while the layer integral of the other term is bounded
by \(C|b'|\eta\), independently of \(L\).  The physical equality
\(b=Lb'\) is not used for that estimate.  In the mixed simple model,
\(\int|\widetilde g|\,\dd t\le C\eta\) and the homogeneous term is
damped for every \(W,L>0\).  The mixed annular fields have a designated
fast factor with a fixed rate margin, giving the bound preceding
\eqref{eq:terminal-signed-slow-left} for every independent bounded
prefactor.  The additive bounds integrate nonconstant powers against
\(\dd w/w\) in (A), and power-corrected fast layers against \(\dd w\)
in (B); they hold uniformly on the closures of the smaller amplitude
ranges and for \(q\ge q_*\).  In the one-rate stable model the fixed
lower large-rate threshold may be increased before defining its domain;
the bounded slow field is then controlled by damping even at that
threshold.  The bounded-\(w_0\) ordinary annular kernel has the guard
in \eqref{eq:terminal-bounded-w-guard}; its forcing and Lipschitz
integrals are at most \(Ch\) on the closed admitted range, even for
independent normalized profile arguments.  Bounded-clock factors have
analytic ODE dependence; if a
finite limiting slow interval is too long for one base-flow chart, take
one fixed finite subdivision and smaller flow charts.  The original
trajectory estimate supplies the same state room at its internal states.
If it is the damping of an old simple kernel that tends to zero, use its
inherited forcing estimate, not a claim about an arbitrary bounded stable
remainder.  In \eqref{eq:terminal-simple-mixed},
\(\int_0^1|L\widetilde g|\,\dd s\le C\eta\) independently of
\(W\ge0\).  In \eqref{eq:terminal-bounded-a},
\[
 \frac{M}{y+\delta}\,\dd s=\frac{\dd y}{y},
\]
so the nonconstant-power estimate bounds its forcing and Lipschitz
integrals independently of \(S\ge0\), including \(S=0\) by continuity.
Thus zero damping is harmless for these particular old kernels; no such
assertion is needed for a general stable equation at zero rate.
Consequently input, output, amplitude, and artificial state-range
inequalities can vanish without loss of these larger existence bounds.

\paragraph{Rate, sign, and unit margins.}
All nonzero rational one-rate or additive rates are fixed positive
numbers.  Mixed (I) has fixed positive layer rates and nonnegative
damping at its limiting boundary.  In the away-from-resonance mixed (II)
rule, let \(r_*>0\) be a lower bound for the finitely many actual
sign-adjusted rates on the chosen coefficient range after shrinking the
generator sizes.  Impose in the independent domain, with slack, the
inequalities \(r_i>r_*/2\).  At their boundary the rates are still
positive.  Here \(r_i=r_i(W/L)\) is the actual sign-adjusted rate of
the independent kernel, not a rate evaluated at a separately copied
ordinary coefficient.  Thus no rate tending to zero at a positive limiting
slope is admitted, even off the physical coefficient graph.  The same choice applies to the
designated fast rates in the small-slope rule.  Its pure-state rates
\(|\beta_i|c\) can tend to zero, precisely when \(c\to0\); the field
then still vanishes without the designated fast layers.  The both-large
case is the zero-slope alternative of the mixed theorem, the bounded
\(W\) case is \eqref{eq:terminal-signed-slow-left}--\eqref{eq:terminal-signed-slow-right},
and bounded clocks are treated below.  Pure slow layers are never used
alone to bound the forcing integral.

Likewise fix the domains of prepared units, root linearizers, and profile
inversions on the larger working boxes where their required Jacobians
and unit divisors are bounded away from zero.  This applies to
\(F,F_s,b,b(0),h\), the denominator in
\eqref{eq:terminal-clock-jacobian}, and the root-chart transversality
constant.  Choose \(|\delta|\le\delta_*<1\) on the working box, so
\(1+\delta/y\) and \(y+\delta\), \(y\ge1\), are units even on the
admitted boundary.  A margin choosing the sign of an algebraic prefactor,
such as \(B\) or \(p\), may instead tend to zero.  The kernel formulas
are analytic in that prefactor at zero; if a normalization divided by it,
that division is an auxiliary graph of the next paragraph, not an
uncontrolled analytic argument.  Overlapping magnitude rules permit the
same conclusion at their finite cutoff boundaries.

\paragraph{Algebraic divisors and positive endpoint variables.}
A physical positive variable or an algebraic normalization divisor may
approach zero.  Keep every bounded normalized value as its own graph
coordinate; for instance \(y=a/d\) uses \(dy=a\), with \(d\ne0\)
retained as a domain condition.  Positive rational powers use polynomial
relations with the designated positive branches.  These relations are
analytic also at zero, although their graph Jacobians may then become
singular.  At each point of the sequence the original nonzero or positive
condition gives the same locally unique graph as before.  A fixed finite
number of differentiations produces only finite powers of these divisors,
which can be cleared without changing the local solution set there.
The coefficients of the resulting equations are analytic at a bounded
graph limit.  If a required normalized value is unbounded, it is instead
a new recoverable large coordinate.  A vanishing determinant imposed by
a later rank restriction behaves in the same way: clear its reciprocal
powers in the finite equations and retain its strict domain condition.
No uniform inverse bound at the limiting boundary point is asserted or
needed for the local graph equivalence at the sequence points.

\paragraph{Zero clocks, ordering boundaries, and ratios.}
The ordinary normalized formulas explicitly include zero length.  The
nozero short field uses its bounded integrated coefficient.  The nozero
long field, if it reaches a bounded-clock corner, retains independent
\(b,b'\) in \eqref{eq:terminal-nozero-extension}.  The simple mixed
field is \(-Wv+L\widetilde g\), whose layers have exponents affine in
\(L\).  Mixed (II) has field \(Lp\mathcal G\) and layers whose
exponents are the fixed affine combinations \(\alpha L+\beta W\).
All are analytic at bounded \((L,W)\), including a zero value of either
clock, without evaluating \(W/L\) in a field coefficient.  A duplicated
bounded coefficient remains independent, with its product tie retained.
If a slope interval itself is imposed using a ratio, its inequalities
and tie may be cleared by \(L>0\); no ratio is necessary inside these
bounded-clock kernels.  For example even \(L,W\to0\) and
\(W/L\to\infty\) cause no singularity in the mixed (I) normalized
field.  A diverging ratio already present as an auxiliary coordinate may
still be retained as a large algebraic input, but does not invoke the
unbalanced mixed primitive when the two primary clocks are bounded.
In particular no auxiliary ratio is introduced merely to express a slope
which these bounded-clock formulas have already eliminated.  The mixed
auxiliary ratio is required only in the specified regimes in which both
primary clocks tend to infinity.

For additive kernels the fixed bound \(q\ge q_*>0\) prevents a zero
denominator.  If \(Q\) stays bounded, then so does
\(M=(Q-q)/q\), and \eqref{eq:terminal-bounded-action} is regular,
including \(Q-q=0\).  For a finite large limit of \(M\), the single
fixed \(K\) in \eqref{eq:terminal-local-ratio} provides analytic
fraction charts.  If \(Q\) is unbounded, the divergent-ratio or
bounded-ratio packet alternatives already proved apply.  If two positive
physical endpoints both approach zero, their ratio is handled by its
graph and \eqref{eq:terminal-log-ties}: bounded log length has the
normalized extension just described, while a divergent log length is an
explicit large coordinate.  Vanishing physical clock speed similarly
affects its algebraic endpoint normalization; the normalized field's unit
denominator remains in its larger working box.

This list exhausts the imposed terminal margins.  Fractions and other
rewritings use an exact identity of the old independent kernel on a
neighborhood of each sequence point, with the same exogenous arguments.
The fixed finite state and jet graphs therefore preserve isolation there,
even if their neighborhoods shrink along the sequence.  If all enlarged
coordinates are bounded, the preceding checks give analytic kernel
factors and analytic divisor-cleared graph equations at their limits.
Otherwise the divergent enlarged coordinates are available for the flag
and packet alternative.  The two conclusions are not conflated with
positive-margin completeness.
\end{proof}

\begin{proof}[Proof of Theorem~\ref{thm:terminal-reduction}]
The real cases above exhaust the constant-state pieces, nozero boxes,
simple boxes, and annuli
from Theorem~\ref{thm:finite-templates}; the nonzero-resonance redo is
finite.  Every additional real cut used a magnitude of a subanalytic
variable, a root or state bin, a coefficient or resonance interval, or a
rational weighting.  Hence these cuts preserve the finite boundary-word
bound.  Boundary connectors keep each actual middle passage compactly
inside its identity piece.  Length alternatives overlap and are selected
from endpoint arguments; large-input subsequences and internal fractions
are used only at an absolute zero test, not as geometric subdivisions.

Choose all analytic boxes first with larger extension, all finite rate
and derivative margins next, and state widths and amplitude ranges finally
inside the required margins.  The nozero integrated prefactors are
explicitly bounded arguments; the simple transfers use damping and small
profile integrals; annular residuals use
\eqref{eq:terminal-profile-integral}.  The additive models have the real
estimates of Theorem~\ref{thm:additive-packets}.  These establish real
existence and state room on open domains described by finitely many strict
inequalities, rather than a condition of unspecified continuation along a
mismatched orbit.  An actual target satisfies tighter margins.  Bounded
tuples whose listed margins stay positive remain in the extended analytic
boxes and have analytic ODE dependence, which proves the completeness
assertion in part (i).  For part (iii), no positive lower bound on these
margins is assumed: Lemma~\ref{lem:terminal-boundary-corners} treats all
their vanishing limits.  Its bounded alternative supplies analytic
normalized kernel factors with state room and divisor-cleared graph
equations; its unbounded alternative supplies the enlarged coordinates
used in the following packet test.

Along an arbitrary test sequence select bounded limits or divergent
subsequences for the finitely many clocks and ratios.  The tables and the
additive ratio alternatives select the asserted primitive.  A finite but
large bounded clock or ratio needs only one fixed finite subdivision for
that subsequence.  Its locally unique state and finite jet graphs preserve
isolation by Lemma~\ref{lem:terminal-ambient-extension}.  The large clocks
have their independent affine flag representation; insert the initial
log nodes and, for an unbalanced mixed test, the auxiliary ratio required
by the corresponding primitive theorem.  On that flag the three packet
theorems supply all fixed finite independent derivatives; bounded
operations and regular pullbacks use Theorem~\ref{thm:calculus}.
No equality between two originally independent clock arguments is imposed
inside a packet test unless provided by its stated flag relation.

These are also the fixed determinations needed on a retest.  The primary
affine clock formulas, lifted logarithms, signs, and analytic profile fields
remain the same under sufficiently small free-coordinate variations
preserving the regime and ordinary limits.  The additive coefficient
determinations are their fixed formal recursions, compatible bounded-charge
coefficients, and, for retained bounded inner endpoints, the normalized
infinite-anchor jets of Theorem~\ref{thm:additive-packets}.  The one-rate
and mixed determinations are the fixed coefficient constructions of their
respective theorems.  Numerical cutoffs or anchors have been retired by
their proved dispersals before they could be discarded from the retained
arguments.  Subsequent regular graph operations therefore act on the same
ordinary coefficient germs by Theorem~\ref{thm:calculus}; inserting unused
nodes uses identity transitions.  This verifies the claimed compatibility
with primitive retests.  Its application to the complete enlarged
matching systems is made in Section~\ref{sec:counting}.

Finally, Lemma~\ref{lem:terminal-tied-endpoints} gives the nearby physical
mismatch identity, and Lemma~\ref{lem:terminal-ambient-extension} gives
the derivative assertion with all graph terms retained.  The common
smallness clauses in the mixed theorem apply because each relevant
residual monomial has a designated uniformly fast layer; hence geometric
tolerances were fixed independently of a subsequently chosen limiting
slope and of all finite asymptotic requests.  This proves every assertion.
\end{proof}
